\documentclass[11pt,a4paper]{article}
\usepackage[T1]{fontenc}
\usepackage[utf8]{inputenc}
\usepackage[english]{babel}
\usepackage{amsmath,amssymb}
\usepackage[margin=2cm]{geometry}
\setlength{\parindent}{0pt}
\setlength{\emergencystretch}{3em}
\allowdisplaybreaks

\title{About \texttt{book\_inequalities\_with\_radicals.tex}}
\author{}
\date{}
\begin{document}
\maketitle
This note explains what the test data in \texttt{book\_inequalities\_with\_radicals.tex} is used for. That file is input for the engine in \texttt{main.c}; it is read by the test \texttt{test\_equations\_inequalities.c} in this folder.

\begin{enumerate}
\item \textbf{Purpose.} The file holds 7 true statements from the same textbook exercises whose printed result contains a square root, such as $-\sqrt{6} < x < 0$. The engine must prove every one of them.

\item \textbf{What it exercises.} Square roots inside a comparison (rule 18p-7): each root becomes a letter $s$ with $s^{2} = R$, and the sign of $P + Q s$ is read from the signs of $P$, $Q$, $R$ and $P^{2} - Q^{2} R$, which are polynomials. Nothing is approximated.

\item \textbf{Variety.} Results with $\sqrt{2}$, $\sqrt{3}$ and $\sqrt{6}$, alone or added to a rational number, in systems and with absolute values.

\item \textbf{Expected result.} The line \texttt{PASS: 7 of 7 statements proved, each with a written proof}.

\item \textbf{How to run it.} From the top folder of the repository:
\begin{verbatim}
cc -std=c11 -O2 -Wall -Wextra -o /tmp/test_equations_inequalities \
    tests/equations_inequalities/test_equations_inequalities.c \
    && /tmp/test_equations_inequalities
\end{verbatim}
\end{enumerate}
\end{document}
